\documentclass[11pt,a4paper]{article}
\usepackage[T1]{fontenc}
\usepackage[utf8]{inputenc}
\usepackage[spanish,es-noquoting,es-noshorthands]{babel}
\usepackage{lmodern,amsmath,amssymb,geometry,hyperref}
\geometry{margin=27mm}
\setlength{\emergencystretch}{3em}
\hypersetup{colorlinks=true,urlcolor=blue,linkcolor=blue}
\DeclareMathOperator{\Spec}{Spec}
\DeclareMathOperator{\Proj}{Proj}
\title{Notas matemáticas sobre EGA II}
\author{}
\date{}
\begin{document}

% BEGIN RIGHTS AND AI AUTHORSHIP NOTICE
\begingroup
\small
\noindent\textbf{Derechos, disponibilidad y autoría}

Los derechos sobre las obras originales y sobre las contribuciones protegidas a estas traducciones y transcripciones derivadas corresponden a sus respectivos titulares. En la medida en que el editor tenga derechos sobre la traducción, la transcripción, las correcciones y la composición producidas por IA, los dedica al dominio público bajo CC0 1.0 Universal (https://creativecommons.org/publicdomain/zero/1.0/deed.es); esto no afecta a los derechos sobre las obras originales ni concede permiso para usarlas. Las menciones de derechos de terceros y las autorizaciones ya concedidas siguen siendo aplicables.

Estas ediciones se ofrecen gratuitamente, con fines no comerciales y exclusivamente educativos, bajo la presunción de que los respectivos titulares de derechos no se opondrían a que estas traducciones y transcripciones estuvieran disponibles gratuitamente junto a los originales franceses en \href{https://www.numdam.org/}{NUMDAM}. Esta presunción no significa que se haya obtenido autorización; el acceso gratuito a los originales no autoriza por sí solo su traducción, adaptación o republicación.

Si usted es titular de derechos sobre algún material presentado aquí y se opone a su disponibilidad, puede \href{https://github.com/KokunoYumeto/ega-fga-es/issues}{abrir una incidencia en GitHub}. El responsable de la publicación retirará el material afectado a petición del titular de los derechos.

Estas notas matemáticas y su composición están producidas por IA. El trabajo documentado incluye OpenAI Codex, GPT-6 Astra, esfuerzo Ultra. Se conserva la autoría de las obras originales y de las contribuciones citadas; no se atribuyen a la IA. El responsable de la publicación no tradujo el texto.
\par
\endgroup
\clearpage
% END RIGHTS AND AI AUTHORSHIP NOTICE

\maketitle
Estas notas acompañan a la edición española con enmiendas de
A.~Grothendieck, con la colaboración de J.~Dieudonné,
\emph{Éléments de géométrie algébrique II},
\emph{Publications Mathématiques de l'IHÉS} \textbf{8} (1961),
pp.~5--222. Los números de resultados sin indicación de volumen
son los de EGA II. La traducción según el impreso y la edición
con enmiendas son textos distintos. Las precisiones matemáticas
de esta última se exponen aquí por separado; no forman parte
del texto de los autores.

La fuente primaria es la edición francesa conservada en
\href{https://numdam.org/item/PMIHES_1961__8__5_0/}{NUMDAM}.
La errata de (3.3.2) citada aquí pertenece a EGA IV-4,
\emph{Publications Mathématiques de l'IHÉS} \textbf{32} (1967),
p.~345. Los referencias a EGA I y EGA III conservan la numeración
de esas obras.

La traducción española utiliza una referencia inglesa y una
estructura TeX derivadas del proyecto humano
\href{https://github.com/ryankeleti/ega/tree/3d7afe28c99d1887ccdb85acf2ca1afd93fb1971}{ryankeleti/ega, versión 3d7afe28}.
Sus colaboradores nombrados son Mateo Carmona, Tim Holzschuh,
Tim Hosgood, Ryan Keleti, Thiago Solovera \& Nery y Wayne Yang;
los \href{https://github.com/ryankeleti/ega/blob/3d7afe28c99d1887ccdb85acf2ca1afd93fb1971/README.md}{créditos por sección}
precisan sus contribuciones. La traducción inglesa heredada
conserva esa autoría. Estas notas y la traducción española no
se atribuyen a esos colaboradores. La atribución no establece
autorización ni aval.


\section{Abiertos principales y secciones sobre una base afín: 1.2.7 y 1.3.7}
Si $S=\Spec(A)$, un abierto principal $D(\lambda)$ se define mediante
un elemento $\lambda\in A$. Para un homomorfismo $\rho:A\to B$,
la expresión $\rho(\lambda)$ pertenece entonces a $B$ y
$f^{-1}(D(\lambda))=D(\rho(\lambda))$. En la demostración de (1.2.7),
el impreso dice $\lambda\in S$; el anillo requerido es $A$.
En (1.3.7), la identidad afín es
$\Gamma(S,\widetilde B)=B$. La letra $U$ que aparece en el impreso
no designa allí un abierto previamente fijado.

\section{Finitud del módulo y del álgebra: 1.4.5}
Las hipótesis de (1.4.3) permiten una $\mathcal O_S$-álgebra
cuasi-coherente arbitraria. Sobre $S=\Spec(A)$ corresponde a una
$A$-álgebra $B$, sin hipótesis de tipo finito sobre $A$.
La equivalencia entre la finitud de $M$ como $B$-módulo y la finitud
de su haz asociado sobre $\Spec(B)$ es I, (1.3.13).
Por ejemplo, $B=A[x_1,x_2,\ldots]$ no tiene por qué ser de tipo finito
sobre $A$, y esa equivalencia sigue siendo válida. La afirmación
adicional de finitud de $B$ y la referencia a I, (9.6.2) en el impreso
no se deducen de las hipótesis ni son necesarias para este argumento.

\section{Propiedad universal del álgebra simétrica: 1.7.4}
La propiedad universal de $\operatorname{Sym}_{\mathcal A}(\mathcal E)$
se formula en la categoría de $\mathcal A$-álgebras conmutativas.
En efecto, un homomorfismo de módulos $u:\mathcal E\to\mathcal B$
se prolonga al álgebra tensorial. Para descender a la simétrica debe
anular las relaciones $x\otimes y-y\otimes x$, es decir,
$u(x)u(y)=u(y)u(x)$. La conmutatividad de $\mathcal B$ asegura esta
condición para todo $u$. Sin ella, dos matrices que no conmuten dan
un homomorfismo desde un módulo libre de rango dos que no admite
esa prolongación. La calificación «conmutativa» explicita la categoría
en que vale la afirmación.

\section{Grados altos y un ideal primo homogéneo: 2.3.2}
Sean $x\in\mathfrak J\cap S_m$ y $f\in S_d\setminus\mathfrak p$,
con $d>0$. La pertenencia $f^rx\in\mathfrak J\cap S_{m+rd}$ vale
para todo $r\geq0$, porque $\mathfrak J$ es un ideal graduado.
La condición que requiere $r$ suficientemente grande es $m+rd\geq n$.
La hipótesis sobre los componentes de grado al menos $n$ da entonces
$f^rx\in\mathfrak p$. Como $\mathfrak p$ es primo y $f\notin\mathfrak p$,
se obtiene $x\in\mathfrak p$; este es el criterio de (2.1.8).
La frase impresa sobre la pertenencia eventual al ideal es verdadera,
pero no separa esa pertenencia de la cota de grado que se utiliza.

\section{Componentes de un homomorfismo graduado: 2.6.4 y 8.14.7}
Si $f\in S_d$ y $\alpha$ es un homomorfismo de grado cero, entonces
$f^n\in S_{nd}$ y $\alpha(f^n)=\alpha_{nd}(f^n)$.
En (2.6.4) las tres componentes aplicadas a las potencias son,
respectivamente, $\alpha_{nd}$, $\alpha_{kd}$ y $\alpha_{nd}$.
En las dos aplicaciones a $f^n$ de (8.14.7) la componente es
$\alpha_{nd}$. Estas graduaciones son también las que permiten
dividir la sección por la potencia invertible sobre $D_+(f)$.
El subíndice $d$ impreso en esos lugares no corresponde al grado
del argumento cuando el exponente varía.

\section{Cartas afines de Proj y componente de grado cero: 3.1.7 y 3.1.8}
En (3.1.7), la identificación del álgebra de la carta $X_f$ con
$\mathcal A$ proporciona la identificación de esquemas mediante
(3.1.4), que describe esa carta como un espectro relativo.
En (3.1.8)(ii), la graduación con componente cero $\mathcal O_Y$
y componentes positivos $\mathcal S_n$ se escribe
\[
\mathcal S'=\mathcal O_Y\oplus\bigoplus_{n>0}\mathcal S_n.
\]
Esto expresa la sustitución de la componente de grado cero utilizada
en la comparación de Proj. La suma con $n\geq0$ que figura en el
impreso incluye además $\mathcal S_0$. No se afirma por ello que
una construcción de álgebra unitaria con esa suma sea imposible;
la lectura indicada precisa cuál es la componente cero de la
construcción considerada.

\section{Homomorfismos de haces asociados: 3.2.6}
Las aplicaciones $\lambda$ y $\mu$ son los homomorfismos canónicos
definidos en (2.5.11.2) y (2.5.12.2). Las hipótesis adicionales de
(3.2.6) son las que permiten afirmar que son isomorfismos.
No puede llamárselas isomorfismos antes de imponer esas hipótesis.
Por ejemplo, para $S=k[t]$ con $\deg t=2$, el haz asociado a $S(1)$
es cero: su localización homogénea tiene componente de grado cero
nula. En cambio, el haz asociado a $S(2)$ es no nulo sobre
$\Proj(S)=\Spec(k)$. Por tanto, la aplicación de producto desde
$\widetilde{S(1)}\otimes\widetilde{S(1)}$ a
$\widetilde{S(2)}$ no es un isomorfismo.

\section{Graduación y localización de un módulo: 3.3.2}
El homomorfismo graduado de (3.3.2.5) tiene por objeto de llegada
$\mathcal M(n)$; la graduación de ese objeto procede de $\mathcal M$.
Esta última lectura figura también en la primera errata para EGA II
de EGA IV-4, p.~345. Sobre una base afín, al localizar se obtiene
\[
M_n\otimes_A S_{(f)}\longrightarrow (M(n))_{(f)},
\qquad x\otimes1\longmapsto x/1.
\]
El objeto $(S(n))_{(f)}$ que aparece en la última línea impresa no es
el módulo localizado de esa aplicación. La graduación y el objeto de
llegada son dos aspectos distintos de la misma construcción.

\section{Denominadores en la generación relativa: 3.4.7}
Para un haz cuasi-coherente sobre $Y=\Spec(A)$, una sección sobre
$D(f)$ se expresa como $s_i/f^n$. Para una familia finita se puede
elegir un mismo exponente $n>0$. Tras el pullback, con $g=\theta(f)$,
se tiene $t_i=s_i g^{-n}$ sobre $p^{-1}(D(f))$.
Así, la igualdad $s|W=\sum_i a_i(t_i|W)$ se convierte en
$s|W=\sum_i b_i(s_i|W)$ con $b_i=a_i(g|W)^{-n}$.
El signo positivo en $s_i g^n$ del impreso no es compatible con
esa fórmula para $b_i$. Por ejemplo, en $k[t,t^{-1}]$, con $n=1$
y $a_i=s_i=1$, los dos signos dan respectivamente $t$ y $t^{-1}$.

\section{Dominio del morfismo de Proj: 3.6.2 y 3.6.4}
Un homomorfismo graduado $\varphi:\mathcal S\to\mathcal S'$
induce un morfismo cuyo dominio de definición $G(\varphi)$ es un
abierto de $\Proj(\mathcal S')$. Cuando $\varphi$ es sobreyectivo,
$G(\varphi)=\Proj(\mathcal S')$ y el morfismo hacia
$\Proj(\mathcal S)$ es una inmersión cerrada. Por eso, en (3.6.2)(i),
el dominio completo es $\Proj(\mathcal S')$, no $X=\Proj(\mathcal S)$.
El resultado que se aplica en (3.6.4) es precisamente (3.6.2)(i);
(3.6.1) es la definición precedente y no tiene un apartado (i).

\section{Dominios locales y secciones que trivializan: 3.8.2 y 3.8.5}
En (3.8.2), $V_\alpha=r^{-1}(D_+(s_\alpha))$ es el abierto sobre
el que se restringe $r$ para obtener el morfismo a $D_+(s_\alpha)$.
La mención a $U_\alpha$ en esa restricción no corresponde al abierto
definido en la frase. En (3.8.5), para cada sección $f_{ij}$ y cada
$x\in X_{f_{ij}}$, es la sección $f_{ij}|V$ la que constituye una
base local del haz invertible sobre un entorno $V$ de $x$.
Los índices vinculan la trivialización con la sección de la familia
finita que se está considerando.

\section{Extensión de escala en la construcción proyectiva: 4.2.1}
En la carta afín de (4.2.1), el objeto de llegada es
$L_r=(S(1))_{(f)}\otimes_{S_{(f)}}B_{[\omega]}$.
La aplicación $E\otimes_A B\to L_r$ viene dada por
\[
x\otimes1\longmapsto (x/1)\otimes1
                 =(f/1)\otimes\omega(x/f).
\]
La igualdad usa el balanceo del producto tensorial sobre $S_{(f)}$.
En particular, el segundo factor del primer tensor es $1\in B$;
el símbolo $f$ impreso en esa posición no expresa esa extensión
de escala.

\section{Haces invertibles en un fibrado proyectivo: 4.2.7}
En la fórmula
$\mathcal L\simeq p^*\mathcal L'\otimes\mathcal O_{\mathbf P(\mathcal E)}(m)$,
la unicidad del entero $m$ requiere fibras proyectivas de dimensión
positiva y una base no vacía. Se consideran, pues, $Y$ no vacío,
conexo y localmente noetheriano y $\mathcal E$ localmente libre de
rango al menos dos. En una fibra, el grado del haz invertible fija
$m$; la conexidad asegura que ese entero localmente constante es
constante. La unicidad de $\mathcal L'$ se deduce de
$p_*\mathcal O_{\mathbf P(\mathcal E)}=\mathcal O_Y$ y la fórmula
de proyección, una vez fijado $m$.

En rango uno, $\mathbf P(\mathcal E)\simeq Y$ y
$\mathcal O(1)\simeq p^*\mathcal E$. Para cada entero $m$ se puede
absorber $\mathcal E^{\otimes m}$ en $\mathcal L'$; por ello $m$
no es único. Sobre la base vacía tampoco se distingue ningún entero.

El cálculo de imágenes directas de III, (2.1.15), junto con la
fórmula de proyección de 0, (5.4.10), da
\[
p_*\mathcal L=0\quad(m<0),\qquad
p_*\mathcal L\simeq\mathcal L'\otimes
                \operatorname{Sym}^m(\mathcal E)\quad(m\geq0).
\]
El haz al que se aplica $p_*$ es $\mathcal L$, no
$\mathcal L^{\otimes m}$: el entero $m$ ya es el exponente de
la torsión en su descomposición. Para describir un morfismo
$\mathbf P(\mathcal E)\to\mathbf P(\mathcal F)$, la aplicación
adyunta $\psi^\sharp:p^*\mathcal F\to\mathcal L$ es un
homomorfismo de $\mathcal O_{\mathbf P(\mathcal E)}$-módulos.
Estas precisiones fijan tanto el exponente como el espacio en que
viven los haces.

\section{Potencias positivas y la inmersión asociada: 4.4.1}
El entero utilizado en (4.4.1)(i) y en su demostración satisface
$n>0$. Se aplica (3.8.5) a una potencia positiva del haz invertible
y se compara la subálgebra generada por las secciones de ese grado
con el correspondiente Veronese. El grado cero no proporciona ese
Veronese ni sustituye a una potencia positiva en la caracterización.
Esta condición no afecta a las sumas de componentes que incluyen
legítimamente la componente de grado cero.

\section{Cuasi-coherencia de tipo finito: 4.4.6 y 4.6.8}
El submódulo $\mathcal E$ de (4.4.6)(b') es cuasi-coherente de
tipo finito. Asimismo, I, (9.4.7), aplicado en (4.6.8), da una
extensión cuasi-coherente de tipo finito del haz de ideales.
No se puede sustituir en general esta conclusión por la coherencia
sin una hipótesis adicional sobre la base. Sobre un anillo no
coherente, incluso un módulo de tipo finito puede no ser coherente.
La finitud de los generadores se mantiene en ambas afirmaciones;
es la propiedad utilizada por los argumentos de inmersión y de
extensión de los ideales.

\section{Clausura y esquema proyectivo: 4.4.7 y 4.5.2}
La existencia de la clausura del subpreesquema que interviene en
(4.4.7) es I, (9.5.10). El apartado siguiente, I, (9.5.11), explica
la terminología de clausura; la referencia a (9.5.10) identifica
directamente la afirmación de existencia empleada.
En la demostración de (4.5.2), $S$ es el anillo graduado de secciones
y el esquema ambiente es $\Proj(S)$. La expresión impresa
$S=\Proj(A)$ confunde ese anillo con el esquema al que se aplica
la inmersión.

\section{Restricción de un haz amplio: 4.5.3}
La definición (4.5.3) de este volumen incluye la cuasi-compacidad
del esquema. Si $\mathcal L$ es amplio y $U\subset X$ es abierto
cuasi-compacto, los abiertos afines de la forma $X_s$ contenidos
en $U$ forman una base para $U$, y las restricciones de las secciones
dan el criterio (4.5.2) para $\mathcal L|U$.
La cuasi-compacidad de $U$ conserva, además, la hipótesis explícita
de la definición. Esta precisión se refiere a la convención del
volumen; no pretende imponer esa misma convención a otras
definiciones de amplitud en la literatura.

\section{Correspondencia de las dos secciones: 4.5.13}
La aplicación de restricción envía la sección $s_\alpha$ de
$\mathcal O_X(p)$ a $s|U_\alpha$, y la sección $t_\alpha$ de
$\mathcal O_X(q)$ a $t|U_\alpha$. Los dos paréntesis introducidos
por «respectivamente» se cierran después de la sección de grado
$q$ y después de su módulo de llegada, respectivamente.
Así quedan separadas las dos correspondencias; no cambia la
identidad de producto de (4.5.13.3).

\section{Abierto de no anulación de una sección inversa: 4.6.13}
La identidad $X_{s'_i}=f^{-1}(Y_{s_i})$ es el apartado (i) del
lema (4.6.13.1), inmediatamente anterior. El localizador
«4.6.1.13» del impreso no corresponde a un resultado de la obra.

\section{Abiertos afines y exactitud de la imagen directa: 5.2.1 y 5.2.2}
En la implicación (c') a (b) de (5.2.1), $U$ es un entorno de $x$.
La sección elegida pertenece al ideal que se anula sobre $X\setminus U$
y satisface $f(x)=1$. En consecuencia,
$x\in X_f\subset U$, no $X_f\subset X\setminus U$.
En (5.2.2), la exactitud es una propiedad del funtor $f_*$ sobre
los haces considerados, cuando varía $\mathcal F$. Para un haz
fijado, $f_*(\mathcal F)$ es un objeto de la categoría de módulos,
no el nombre del funtor cuya exactitud se invoca.

\section{El grafo en la demostración del lema de Chow: 5.6.1}
En la parte (D), los abiertos previamente definidos son $U''_k$.
El grafo $T_k\subset X\times_S W_k$ de $h_k:W_k\to X$ es cerrado,
pues $X$ es separado sobre $S$, y su proyección a $W_k$ es un
isomorfismo. El morfismo desde el abierto $U'$ se factoriza a través
de ese grafo por la conmutatividad de (5.6.1.5).
La propiedad de clausura de I, (9.5.8), extiende la factorización
al subpreesquema inducido por $X'$ sobre $U''_k$. La proyección del
grafo da entonces la inmersión requerida de $U''_k$ en $W_k$.
Hablar de esta factorización precisa los morfismos a los que se
aplica la propiedad de clausura.

\section{Composición con un grafo cerrado: 5.6.3}
En la demostración, $X$ es propio sobre $S$, $P$ es separado sobre
$S$, $f:X'\to X$ es proyectivo y $j:X'\to P$ es una inmersión.
La factorización con todos sus dominios es
\[
X'\xrightarrow{\Gamma_j}X'\times_S P
\xrightarrow{f\times1_P}X\times_S P
\xrightarrow{q_2}P.
\]
Su compuesto es $j$. El grafo $\Gamma_j$ es cerrado por la
separación de $P$. El segundo morfismo es propio por cambio de
base de $f$, y $q_2$ es propio por cambio de base de $X\to S$.
Por ello $j$ es cerrado; al ser ya una inmersión, es una inmersión
cerrada. La expresión que compone solo $f\times1_P$ y $q_2$
omite el primer morfismo y tiene por dominio $X'\times_S P$,
no $X'$.

\section{Morfismos finitos, fibras y productos: 6.1.10--6.2.4}
El paso a los esquemas reducidos en (6.1.10) utiliza (6.1.5)(vi);
esa proposición consta de los apartados (i)--(vi), sin apartado (vii).
La finitud de las fibras que permite concluir en (6.2.3) que un
morfismo finito es cuasifinito es (6.1.7), junto con el criterio
de (6.2.2). El resultado (6.1.8) trata otra propiedad de los
cuerpos de funciones.
Para $f:X\to Y$ y $g:X'\to Y'$, el producto de (6.2.4)(iv) es
\[
f\times_S g:X\times_S X'\longrightarrow Y\times_S Y'.
\]
Los dos factores del dominio son las fuentes de $f$ y $g$;
los del codominio son sus objetos de llegada.

\section{Rango y número de generadores: 6.4.3}
El entero $n$ de la demostración es el rango de $E$, reducido a uno
mediante la potencia exterior. Ese rango no es el número de
generadores de $E$ sobre $A$. Si $e_1,\ldots,e_m$ generan $E$,
la aplicación
\[
\operatorname{Hom}_A(E,E)\longrightarrow E^m,
\qquad u\longmapsto(u(e_1),\ldots,u(e_m))
\]
es inyectiva: un endomorfismo nulo sobre los generadores es nulo.
Se requiere, por tanto, un índice independiente $m$. Por ejemplo,
el ideal $(x,y)$ de $k[x,y]$ tiene rango uno, pero no es cíclico.

\section{Cuerpos de fracciones en el criterio valuativo: 7.1.4}
La condición (i) exige un $\kappa(y)$-homomorfismo
$\kappa(x)\to\kappa(b)$, donde $b$ es el punto genérico de
$Y'=\Spec(A')$. El cuerpo $\kappa(b)$ es $\operatorname{Frac}(A')$;
no es el cuerpo residual del punto cerrado de $A'$.
Por (7.1.2), se elige primero un anillo de valuación $B$ que domine
$\mathcal O_{y'}$ y cuyo cuerpo de fracciones sea $\kappa(x)$.
Después se elige un anillo de valuación completo $A'$ que domine
$B$ y tenga cuerpo residual algebraicamente cerrado. La inclusión
$B\subset A'$ induce
\[
\kappa(x)=\operatorname{Frac}(B)
\longrightarrow\operatorname{Frac}(A')=\kappa(b).
\]
Para la condición (ii), basta tomar directamente el anillo de
valuación que domina $\mathcal O_{y'}$ con cuerpo de fracciones
$\kappa(x)$. Una condición sobre la extensión del cuerpo residual
de $A'$ no proporciona por sí sola el homomorfismo entre cuerpos
de fracciones que exige (i).

\section{Imagen de una especialización: 7.2.1}
En (7.2.1)(b), $y'$ es una especialización de $y=f(x)$.
El punto buscado $x'$ es una especialización de $x$ situada sobre
$y'$, de modo que la igualdad es $f(x')=y'$.
La letra $x$ que figura en el segundo miembro del impreso pertenece
al espacio fuente y no expresa la condición de estar sobre $y'$.

\section{El radical y la graduación natural: 8.2.10}
Con $S=\bigoplus_{d\geq0}S_d$, se escribe
$S_{[n]}=\bigoplus_{d\geq n}S_d$ y
$S^\natural=\bigoplus_{n\geq0}S_{[n]}$, graduado por $n$.
Para una familia de elementos homogéneos positivos $(f_\alpha)$,
la condición pertinente es
\[
S_+\subset\sqrt{(f_\alpha)S},
\quad\text{y la conclusión es}\quad
S^\natural_+\subset\sqrt{(f^\natural_\alpha)S^\natural}.
\]
Los radicales se toman en los anillos unitarios indicados. No se
exige que coincidan con todo el ideal irrelevante: pueden contener
nilpotentes de grado cero.

Para comprobar la conclusión, considérese $g\in S_{k+h}$ como
elemento de grado $k>0$ de $S^\natural$, con $h\geq0$.
La hipótesis da, para algún exponente y después para todos los
exponentes mayores,
$g^m=\sum_\alpha c_{\alpha m}f_\alpha$, con un conjunto finito
fijo de índices. Tomando componentes homogéneas, los coeficientes
tienen grado $m(k+h)-\deg f_\alpha$ en $S$.
Para $m$ suficientemente grande, $km>\deg f_\alpha$ para todos
esos índices. El mismo coeficiente se puede considerar como
elemento de grado $km-\deg f_\alpha>0$ de $S^\natural$, pues
\[
m(k+h)-\deg f_\alpha\ \geq\ km-\deg f_\alpha.
\]
La igualdad expresa entonces $g^m$ en el ideal generado por los
$f^\natural_\alpha$. Los elementos homogéneos positivos así
considerados generan $S^\natural_+$ como grupo aditivo; el radical
es un ideal, de modo que se obtiene la inclusión anunciada.

Para ver por qué no se debe exigir igualdad con el radical de todo
el anillo, tómese $S=k[\varepsilon]/(\varepsilon^2)$ concentrado
en grado cero y la familia vacía. Se tiene $S_+=0$, mientras que
$\sqrt{(0)}=(\varepsilon)\ne0$. La formulación mediante inclusión
permite esta situación y expresa el argumento de (8.2.10)(iv).

\section{El álgebra del lema de generación: 8.2.14.1}
El lema tiene por álgebra $T$ y su conclusión compara $T_n$ con
$T_1^n$ para $n$ suficientemente grande. La descomposición en
monomios de sus generadores $h_j$ se aplica, por tanto, a un
elemento de $T_n$. La letra $S_n$ de esa frase pertenece a la
notación de la proposición circundante, no a la del lema.

\section{Soporte finito local de una acción de grupo multiplicativo: 8.3.9}
Una aplicación de haces
$\psi:\mathcal S\to\mathcal S[T,T^{-1}]$ se escribe
$\psi=\sum_{n\in\mathbf Z}\psi_nT^n$.
La suma directa que define el haz de polinomios de Laurent tiene
soporte finito local: para una sección $s$ y un punto, existe un
entorno de ese punto en el que solo un número finito de
$\psi_n(s)$ son no nulos. En un abierto arbitrario no tiene por
qué existir un único conjunto finito de índices válido en todo él.

Por ejemplo, sea $Y=\coprod_{r\geq1}\Spec(k)$ y, en cada componente,
considérese $k[z]$ con la acción $z\mapsto zT^r$.
La familia de secciones $z$ define una sección sobre $Y$.
En cada componente su imagen tiene un solo peso, pero al recorrer
$Y$ aparecen todos los pesos $r$. Esto demuestra la necesidad de
la calificación local. Sobre un abierto cuasi-compacto, un
recubrimiento finito por tales entornos sí proporciona un conjunto
finito común de índices.

\section{La definición de la graduación natural: 8.7.3}
La igualdad $S_n^\natural=S_{[n]}$ utilizada en la demostración
es la definición (8.2.6.2). El número «8.2.7.4» del impreso no
corresponde a una fórmula de (8.2.7), que termina en (8.2.7.3).
La referencia a (8.2.6.2) fija la construcción a la que se aplica
la localización de la demostración.

\section{Contravarianza y componentes de un álgebra: 8.7.4}
Un homomorfismo $\mathcal S''\to\mathcal S$ induce por
contravarianza un morfismo
$\Spec(\mathcal S)\to\Spec(\mathcal S'')$.
Lo mismo ocurre con el morfismo de Proj en su dominio de definición.
Así, el morfismo vertical de (8.7.4)(i) va de $C_X$ a $C_{X''}$,
y el de la derecha va de $\Proj(\mathcal S^\natural)$ a
$\Proj(\mathcal S^{\prime\prime\natural})$.
Estas direcciones hacen que las flechas horizontales canónicas
formen un cuadrado con los mismos dominios y codominios.
En (8.7.4)(ii), el álgebra graduada es
$\mathcal T=\bigoplus_{n\geq0}\mathcal J_n$; el componente del
sumando de índice $n$ es $\mathcal J_n$, no $\mathcal J_m$ con
un índice libre.

\section{La variable formal y las componentes de la fibra: 8.10.1 y 8.10.3}
El isomorfismo de (8.10.1)(iii) tiene por dominio $A[[T]]$ y por
codominio $C=\varprojlim B/\mathfrak J^{n+1}$.
Envía la variable formal $T$ al elemento $t\in C$; esta dirección
determina la aplicación sobre las series mediante completitud.

En (8.10.3), el anillo local regular de dimensión uno es el de la
fibra $\pi^{-1}(x)$ en $j(x)$. La reducción a
$\Spec(\kappa(x))$ convierte el espacio total del argumento en esa
fibra. Su anillo local es íntegro, de modo que existe una única
componente irreducible de la fibra que contiene $j(x)$, por I,
(5.1.4). Escribir explícitamente $\pi^{-1}(x)$ conserva el objeto
del enunciado; la $V$ del argumento reducido designa ese cambio
de base, no una afirmación nueva sobre las componentes del espacio
total anterior al cambio de base.

\section{Hipótesis del morfismo de contracción: 8.11.7}
La demostración invoca (8.11.6). Además de las hipótesis impresas
de (8.11.7), se requieren las dos condiciones sobre $q$ de ese
resultado: $q$ es propio y
\[
q:V\setminus j(X)\ \xrightarrow{\ \sim\ }\
       C\setminus\varepsilon(Y).
\]
Con ellas, (8.11.6) proporciona la biyectividad y el carácter
cerrado usados para $u:C'\to C$; la identificación fuera de la
sección proporciona la birracionalidad. El resultado (8.11.2),
aplicado al destino normal, da entonces el isomorfismo.

Sin esas condiciones, tómense $Y=C=\Spec(k)$,
$X=\mathbf P^1_k$, $\mathcal L=\mathcal O(1)$ y el morfismo
estructural $q:\mathbf V(\mathcal L)\to\Spec(k)$.
Las condiciones impresas sobre $p$, la amplitud relativa, la
normalidad de $C$ y $p_*\mathcal O_X=\mathcal O_Y$ se cumplen,
y el diagrama con las secciones es conmutativo. Sin embargo,
\[
C'=\Spec\!\left(\bigoplus_{n\geq0}
              H^0(\mathbf P^1_k,\mathcal O(n))\right)
   =\Spec(k[x_0,x_1])=\mathbf A^2_k,
\]
que no es isomorfo a $C$. Las condiciones adicionales excluyen
este ejemplo y permiten la aplicación citada en la prueba.

\section{Identificación y restricción de un haz: 8.12.8 y 8.13.3}
La reducción de (8.12.8) usa los isomorfismos (8.12.5.3) y
(8.12.6.1); la primera referencia identifica la comparación de
haces necesaria para el homomorfismo canónico considerado.
En (8.13.3), $\mathcal N$ es el módulo sobre la base, mientras
que $\widetilde{\mathcal N}$ es el haz sobre el espectro relativo
$C$. El objeto que se restringe a
$C\cap\widehat C_f=\Spec(S_f)$ es
$\widetilde{\mathcal N}$, conforme a (8.3.2.6).

\section{Estructura de módulo de las secciones graduadas: 8.14.5}
El haz $\mathcal F$ de (8.14.5) es un módulo. La multiplicación
por secciones de $\mathcal S_X$ induce en
$\boldsymbol\Gamma_*(\mathcal F)$ una estructura de módulo
graduado sobre $\boldsymbol\Gamma_*(\mathcal S_X)$.
No hay en las hipótesis una multiplicación interna ni una unidad
en $\mathcal F$ que lo conviertan en un álgebra. La estructura
afirmada es, por tanto, la de módulo graduado.

\section{Notaciones y bibliografía}
En el índice de notaciones, la norma de (6.5.1) se aplica a una
sección $f$ de la $\mathcal A$-álgebra $\mathcal B$.
El cono proyectante se define en (8.3.1), que es el localizador
correspondiente a esa entrada.
La bibliografía impresa de EGA II es la continuación [23]--[26]
en la página 205. Las entradas [1], [9], [11] y [13] que se citan
en el volumen pertenecen a la bibliografía precedente de la serie.
Su presencia en el lector autónomo con enmiendas permite consultar
esas referencias sin atribuirlas a la continuación impresa de EGA II.

\section{Objetos representados en los addenda a EGA I}
En la demostración de I, (1.8.3), incluida en los addenda de EGA II,
los funtores comparados son $\operatorname{Hom}(-,Y)$ y
$\operatorname{Hom}(-,\Spec(A))$. Su isomorfismo natural da,
por la propiedad de los funtores representables, un isomorfismo
$Y\to\Spec(A)$; la letra $X$ del impreso no es el objeto representado.
La aplicación sobre pares provistos de un módulo que se usa en
I, (1.8.10), es (1.8.9.1), definida inmediatamente antes.
En I, (1.8.11), el resultado (1.8.7) es un corolario; su
identificación y su aplicación a los funtores considerados no
dependen de llamarlo proposición.

\end{document}
